\emph{Preliminaries.} By Hypothesis (S), clause (v) (\claimref{standard-model.hyp-s-slot-statistics}), the five modes are one Clifford system, and the pair-number form of $B - L$ is that of Proposition \claimref{standard-model.under-s-composites-state-space-spinor}(3). $v_{i}^{2} = w_{i}^{2} = 0$ and $v_{i}w_{j} + w_{j}v_{i} = \delta_{ij}$, distinct modes anticommuting, so $n_{i}^{2} = v_{i}(1 - v_{i}w_{i})w_{i} = n_{i}$ and the $n_{i}$ commute. $1 - 2n_{i} = 1 - \tfrac12(e_{i} + f_{i})(e_{i} - f_{i}) = e_{i}f_{i}$, so $n_{i} - \tfrac12 = -\tfrac12 e_{i}f_{i}$. $[n_{i}, v_{i}] = v_{i}w_{i}v_{i} = v_{i}(1 - v_{i}w_{i}) = v_{i}$ and likewise $[n_{i}, w_{i}] = -w_{i}$, while $n_{i}$ commutes with the vectors of the other modes. Hence the $4^{5}$ ordered monomials $v_{S}w_{S'}$ (a basis of $\Cl(5,5)$) are eigenvectors of $x \mapsto [k, x]$ with eigenvalue $d = |S \cap \{1,2,3\}| - |S' \cap \{1,2,3\}|$, and $\Cl(5,5) = \bigoplus_{d = -3}^{3}\mathcal{A}_{d}$. The element $k$ satisfies $k(k-1)(k-2)(k-3) = 0$ (it is a sum of three commuting idempotents), so $1 = p_{0} + p_{1} + p_{2} + p_{3}$ with $p_{\kappa}$ the Lagrange polynomials in $k$, orthogonal idempotents with $k\,p_{\kappa} = \kappa\,p_{\kappa}$. On the range of $p_{\kappa}$ an odd number of the factors $1 - 2n_{a}$ equal $-1$ exactly when $\kappa$ is odd, so $P\,p_{\kappa} = (-1)^{\kappa}p_{\kappa}$.

\emph{(i)} By the last remark, $n_{B}\,p_{\kappa} = \tfrac14\bigl(1 - \tfrac23\kappa - (-1)^{\kappa}\bigr)p_{\kappa}$, which is $0, \tfrac13, -\tfrac13, 0$ for $\kappa = 0, \dots, 3$, and $n_{L}\,p_{\kappa} = -\tfrac14\bigl((-1)^{\kappa} + 3 - 2\kappa\bigr)p_{\kappa}$, which is $-1, 0, 0, 1$. Since $(1 - 2n_{s})^{2} = 1$, $\chi\Gamma = -(1 - 2n_{4})^{2}(1 - 2n_{5})^{2}P = -P$; where $\Gamma = 1$ this gives $-P = \chi$ and the definitions become the stored ones, whose values on the species of the sixteen are those of Proposition \claimref{standard-model.proton-decay-only-through-slot-pair}; where $\Gamma = -1$, $-P = -\chi$. The complement of a state has pair number $3 - \kappa$, and both lists of values are odd under $\kappa \mapsto 3 - \kappa$. From the preliminaries, $k - \tfrac32 = -\tfrac12\sum_{a}e_{a}f_{a}$, so $B - L = -\tfrac23(k - \tfrac32) = \tfrac13\sum_{a}e_{a}f_{a}$, and $P = e_{1}f_{1}e_{2}f_{2}e_{3}f_{3} = -e_{1}e_{2}e_{3}f_{1}f_{2}f_{3}$ (three transpositions bring the $f$ to the right), which gives the displayed form of $n_{B}$. Finally $e_{1}$ anticommutes with $e_{1}f_{1}$ and commutes with $e_{2}f_{2}$ and $e_{3}f_{3}$, so $[P, e_{1}] = 2Pe_{1}$, of grade five and not zero, and $[n_{B}, e_{1}] = -\tfrac16 f_{1} - \tfrac12 Pe_{1}$; the same holds for $e_{2}, e_{3}$. The elements whose commutator preserves the space of vectors and generates a rotation are the bivectors, and $n_{B}$ is not of that kind.

\emph{(ii)} Let $g$ be a polynomial and $x \in \mathcal{A}_{d}$. Then $kx = x(k + d)$, so $g(k)\,x = x\,g(k + d)$ and
\[
[g(k), x] = x\bigl(g(k + d) - g(k)\bigr) = \sum_{\kappa}\bigl(g(\kappa + d) - g(\kappa)\bigr)\,x\,p_{\kappa},
\]
where $x\,p_{\kappa} = 0$ unless $\kappa + d \in \{0, 1, 2, 3\}$ (because $k\,x\,p_{\kappa} = (\kappa + d)\,x\,p_{\kappa}$ and $k$ has no other eigenvalue). The elements $x\,p_{\kappa}$ for different $\kappa$ are linearly independent when nonzero, the $p_{\kappa}$ being orthogonal idempotents. So $[g(k), x] = 0$ exactly when $x\,p_{\kappa} = 0$ for every $\kappa$ with $g(\kappa + d) \ne g(\kappa)$. Take $g$ with $g(k) = n_{B}$, values $0, \tfrac13, -\tfrac13, 0$. For $d = 0$ no condition arises. For $d = \pm 3$ the only admissible step is between $\kappa = 0$ and $3$, where the values agree. For $d = \pm 1$ the steps have differences $\tfrac13, -\tfrac23, \tfrac13$, all nonzero, and for $d = \pm 2$ they have $-\tfrac13, -\tfrac13$; so $x\,p_{\kappa} = 0$ for all $\kappa$ and $x = \sum_{\kappa}x\,p_{\kappa} = 0$. Since $n_{B}$ commutes with $k$, the commutator with $n_{B}$ preserves each $\mathcal{A}_{d}$, and an element commutes with $n_{B}$ exactly when each of its components does. This proves the first statement. For $B - L$, $g$ is linear with $g(\kappa + d) - g(\kappa) = -\tfrac23 d$, nonzero for $d \ne 0$. The dimensions count monomials: on the three pairs there are $\sum_{r}\binom{3}{r}^{2} = 20$ with $d = 0$ and one each with $d = \pm 3$, and $16$ monomials in the slot vectors.

\emph{(iii)} The bivectors are spanned by the products of two distinct vectors among the $v_{i}, w_{i}$ and the $n_{i} - \tfrac12$, each homogeneous: $d = 0$ for the nine $v_{a}w_{b}$, $n_{a} - \tfrac12$ and for the six slot bivectors, $d = \pm 2$ for the six $v_{a}v_{b}$, $w_{a}w_{b}$, $d = \pm 1$ for the twenty-four products of a slot vector with a pair vector. No bivector has $d = \pm 3$, so by (ii) the bivectors commuting with $n_{B}$ are the fifteen; the six slot bivectors span $\su(2)_{L} \oplus \su(2)_{R}$ and the fifteen pair bivectors $\su(4)$ by Proposition \claimref{standard-model.automorphisms-preserve-slot-pair-split-are}. $P$ anticommutes with each pair vector and commutes with each slot vector, and $\chi$ the reverse; so $P$ and $\chi$ commute with the six and anticommute with the twenty-four. For $x$ among the six, $[n_{B}, x] = \tfrac14[B - L, x] - \tfrac14[P, x] = \tfrac14 \cdot (-\tfrac23 d)\,x$ with $d = \pm 2$. For $x$ among the twenty-four, say $d = 1$, the formula of (ii) gives $[n_{B}, x] = \tfrac13 x p_{0} - \tfrac23 x p_{1} + \tfrac13 x p_{2}$; $x = v_{s}v_{a}$ or $w_{s}v_{a}$ does not annihilate the states in which pair $a$ is empty, which occur at $\kappa = 0$ and at $\kappa = 1$, so $x p_{0}$ and $x p_{1}$ are both nonzero and independent, and the combination is not a multiple of $x = x p_{0} + x p_{1} + x p_{2}$. The case $d = -1$ is the same with pair $a$ occupied, at $\kappa = 1$ and $\kappa = 2$.

\emph{(iv)} $T$ is the Majorana element of Proposition \claimref{standard-model.majorana-element} and the rung vectors are those of Proposition \claimref{standard-model.bidoublet-complex}. $T \in \mathcal{A}_{3}$, so $[n_{B}, T] = 0$ by (ii) and $[B - L, T] = -\tfrac23 \cdot 3\,T = -2T$, as stored. $(1 - 2n_{a})v_{a} = -v_{a}$ and $v_{a}(1 - 2n_{a}) = v_{a}$, so $PT = -T$, $TP = T$ and $[P, T] = -2T$; hence $[n_{L}, T] = -\tfrac14(-2T) - \tfrac34(-2T) = 2T$. $T^{\dagger} = w_{5}w_{4}w_{3}w_{2}w_{1} \in \mathcal{A}_{-3}$. A slot vector commutes with $k$, with $P$ and so with all three charges. A pair vector $v_{a}$ lies in $\mathcal{A}_{1}$ and is not zero.

\emph{(v)} The linear map fixing the $e_{i}$ and negating the $f_{i}$ preserves $Q$ and so extends to an automorphism $X$. With $u = e_{1}\cdots e_{5}$, $u^{2} = (-1)^{10} = 1$, $u\,e_{i}\,u = (-1)^{4}e_{i}$ and $u\,f_{i}\,u = (-1)^{5}f_{i}$, so $X$ is conjugation by $u$. $X(v_{i}) = w_{i}$, $X(n_{i}) = w_{i}v_{i} = 1 - n_{i}$, so $X(k) = 3 - k$, $X(B - L) = -(B - L)$, $X(P) = (-1)^{3}P$, $X(\Gamma) = (-1)^{5}\Gamma$, $X(\chi) = \chi$, and $X(n_{B})$, $X(n_{L})$ follow. $X(T) = w_{1}\cdots w_{5} = (-1)^{10}w_{5}\cdots w_{1} = T^{\dagger}$. $X$ maps $\mathcal{A}_{d}$ to $\mathcal{A}_{-d}$ and preserves the split of the vectors into slot and pair vectors, so it maps each family of (iii) onto itself; $\su(2)_{L}$, spanned by $v_{4}w_{5}$, $v_{5}w_{4}$, $\tfrac12(n_{4} - n_{5})$, goes to the span of $w_{4}v_{5} = -v_{5}w_{4}$, $w_{5}v_{4} = -v_{4}w_{5}$, $-\tfrac12(n_{4} - n_{5})$.

\emph{(vi)} The centraliser of $n_{B}$ is a subalgebra and by (ii) to (iv) contains every generator of $\mathcal{D}$; the argument of (ii) is unchanged by an extension of the scalars. $T \in \mathcal{D}$ commutes with neither $B - L$ nor $n_{L}$ by (iv). The last sentence is (iii) and (iv).
